\section{Iterative Baseline Extension}

\subsection{Motivation and Objective}

Building on the unified execution pipeline introduced in previous phases, this task
extends both the SQL and NL baselines with a controlled iterative retrieval mechanism.
The goal is to improve result completeness by allowing the LLM to generate multiple
partial result sets across successive calls, while preserving full compatibility with the
project’s standardized JSON output schema.

This mechanism addresses cases where a single LLM invocation fails to recover the
entire expected answer set due to truncation, prompt limitations, or generation
variability. Iteration provides a systematic way to accumulate non-overlapping tuples
across rounds.

\subsection{Design Principles and Architectural Integration}

The iterative enhancement was introduced without modifying the core architecture.
A configurable parameter \texttt{max\_iter} (defined in \texttt{config.yaml}) controls the
behaviour and ensures consistent execution across providers.

The integration adheres to the following principles:

\begin{itemize}
    \item \textbf{Provider-Agnostic Execution.} Iteration is implemented at the baseline level 
    (SQL and NL), leaving LLM instantiation and LCEL chain construction unchanged.
    This guarantees full compatibility with any provider supported through the Adapter 
    and Factory patterns.

    \item \textbf{Minimal Invasiveness.} No modifications were made to the LCEL chain,
    output parser, or prompt template. Iteration is strictly an additional layer wrapped
    around the existing call--parse--store cycle.

    \item \textbf{Deduplicated Accumulation.} Each iteration contributes new tuples 
    based on strict key-wise comparison. A local accumulator collects all rows returned 
    in previous rounds, preventing redundancy and guaranteeing stable output.

    \item \textbf{Early Stopping Criterion.} To avoid unnecessary LLM calls, the system
    terminates early if an iteration yields zero new rows. This rule applies identically
    to SQL and NL baselines.
\end{itemize}

\subsection{Iterative Execution Workflow}

For each query (SQL baseline) or natural language prompt (NL baseline), the extended
workflow proceeds as follows:

\begin{enumerate}
    \item \textbf{Initialization.}  
    Create an accumulator for unique tuples and load \texttt{max\_iter}.

    \item \textbf{Repeated Invocation.}
    At each iteration, the baseline executes:
    \begin{enumerate}
        \item construction of the LCEL chain with the existing prompt template,
        \item LLM invocation using the configured provider,
        \item parse the response via the Pydantic output parser.
    \end{enumerate}

    \item \textbf{Deduplication and Union.}  
    Compare new tuples to the accumulator using canonical JSON formatting; retain
    only unseen tuples.

    \item \textbf{Stopping Condition.}  
    Terminate when no new tuples are found.

    \item \textbf{Persistence.}  
    After early stopping or reaching \texttt{max\_iter}, serialize:
    \begin{itemize}
        \item the union of all retrieved tuples,
        \item total elapsed time,
        \item aggregated token usage.
    \end{itemize}
\end{enumerate}

\subsection{Baseline-Specific Considerations}

\paragraph{SQL Baseline.}  
Iteration compensates for incomplete tuple enumeration in a single pass. Because SQL
inputs are fully specified, repeated calls primarily improve recall.

\paragraph{NL Baseline.}  
Iteration mitigates incomplete enumerations and semantic drift in free-form queries.
Early stopping is essential to prevent hallucinated or reformulated rows once the LLM
stabilizes on a fixed subset.

\subsection{Outcome and Observations}

Controlled iteration increases tuple coverage while remaining fully aligned with the
evaluation framework. The design preserves modularity, provider-agnosticism, and the
standardized JSON format established in previous phases.

The mechanism also maintains conceptual continuity with iterative operators in the
GALOIS framework (e.g., Key-Scan), providing a foundation for more advanced
retrieval strategies in future work.

\section{Evaluation of the Iterative Baselines}
We evaluated the iterative SQL and NL baselines on the full benchmark suite
(FLIGHT-2, FLIGHT-4, GEO, MOVIES, PRESIDENTS, WORLD).

\subsection{Iterative Execution Results}

The following tables report iterative results for both LLM providers.

\begin{table}[h!]
\centering
\begin{tabular}{lcc}
\toprule
Metric & NL & SQL \\
\midrule
F1-CELL       & 0.425 & 0.456 \\
CARDINALITY   & 0.696 & 0.607 \\
TUPLE-CONSTR. & 0.262 & 0.245 \\
\#TOKENS (M)   & 0.028 & 0.067 \\
AVG-TIME (s)  & 5.035 & 9.681 \\
\textbf{AVG-SCORE} & \textbf{0.461} & \textbf{0.436} \\
\bottomrule
\end{tabular}
\caption{Iterative Results for \texttt{gemini-2.5-flash}.}
\label{tab:iter_gemini}
\end{table}

\begin{table}[h!]
\centering
\begin{tabular}{lcc}
\toprule
Metric & NL & SQL \\
\midrule
F1-CELL       & 0.271 & 0.370 \\
CARDINALITY   & 0.500 & 0.451 \\
TUPLE-CONSTR. & 0.091 & 0.183 \\
\#TOKENS (M)   & 0.005 & 0.006 \\
AVG-TIME (s)  & 2.206 & 2.812 \\
\textbf{AVG-SCORE} & \textbf{0.287} & \textbf{0.335} \\
\bottomrule
\end{tabular}
\caption{Iterative Results for \texttt{meta-llama/llama-3-70b}.}
\label{tab:iter_llama70b}
\end{table}

The iterative execution impacts the behaviour of both baselines in distinct
ways. In the SQL setting, the model benefits from task repetition, often
recovering additional valid tuples that were missed in the first pass. This
results in higher tuple–level completeness and, in many cases, improvements in
the \textsc{F1-Cell} and \textsc{Cardinality} metrics.  
For NL queries, which are inherently more open–ended, the iterative mechanism
still enhances tuple coverage, though with a less uniform effect across
datasets.

\subsection{Comparison with Non-Iterative Baselines}
To contextualise the effect of iteration, we compare the results above with the
single-pass baselines. Overall, the iterative variant provides measurable gains 
in tuple completeness and global \textsc{Avg-Score}, especially in datasets where 
the model tended to omit large fractions of the answer in a single call.  
As expected, the iterative mechanism also increases the total number of
generated tokens and the average execution time. However, these costs remain
moderate and proportional to the configured value of \texttt{max\_iter}.

In conclusion, iterative execution improves robustness and recall, providing a
more stable foundation for evaluating LLM behaviour under both SQL and NL
interfaces.


